\section{Running the analysis}
\label{sec:running}

The top of the right sidebar holds the four buttons that start, stop and use a run, and the
\ui{PROGRESS} area below them (Figure~\ref{fig:progress}). This section follows a run from the
click on \ui{Run DIC Analysis} to the moment its results appear.

\screenshot[0.42\textwidth]{fig-progress.png}{The top of the right sidebar during a run:
\ui{Run DIC Analysis} (disabled while the run goes on), \ui{Cancel}, \ui{Export Results},
\ui{Open Strain Window} and the \ui{PROGRESS} area with its bar, the status line and the
\ui{ELAPSED} and \ui{REMAINING} times.}{fig:progress}

\subsection{Run DIC Analysis}

\ui{Run DIC Analysis} tracks every frame of the sequence with the tracking mode, the model and
the parameters set in the left sidebar (Sections~\ref{sec:workflow} and~\ref{sec:parameters}).
The button is enabled whenever no run is going on. Before anything is computed, RAFTcorr checks
the inputs; a missing input opens a \ui{Cannot start analysis} dialog that says what to do:

\begin{itemize}
  \item fewer than 2 images: \emph{Need at least 2 images. Load an image folder in the IMAGES
        section before running DIC.}
  \item no region of interest on frame~1: \emph{Define a Region of Interest first: draw or
        import one for frame 1 in the REGION OF INTEREST section.}
  \item a region of interest on frame~1 with no pixel selected, such as an all-black mask from
        \ui{Batch Import}: \emph{The Region of Interest of frame 1 is empty (no pixels selected).
        Draw it again in the REGION OF INTEREST section.}
  \item refused text in the \ui{Reference Frames} field, with \ui{Custom Frames}: the dialog
        repeats the reason the field gives (Section~\ref{sec:workflow});
  \item a setting the run cannot use: the same dialog names it.
\end{itemize}

In incremental tracking, a reference frame without a region of interest of its own does not stop
the run: RAFTcorr carries the region of frame~1 to it, and the \ui{LOG} lists such frames
(Section~\ref{sec:workflow}).

The run then starts in the background. The \ui{LOG} records what it does: the model, the number
of frames, the tracking mode, the number of iterations and the device; the line \emph{loaded
RAFTcorr-s4 (\ldots)} when the model is loaded into memory (it stays loaded for the next run with
the same model and device); the size and position of the region of interest; how the region is
cut into windows (Section~\ref{sec:tiling}). A frame whose tracked part covers less
than 90\,\% of the region of interest gets a warning line, \emph{Frame $k$: only \ldots{} of the
region of interest could be tracked}.

Results appear only when the run ends: there is no preview of frames during the run. Until then
the canvas keeps showing the previous results, if any. They are replaced only when the new run
returns at least one computed frame. Results larger than 2\,GB are kept on disk in a temporary
folder until you quit (\emph{results of \ldots{} GB are kept on disk in \ldots},
Section~\ref{sec:overview-files}).

Loading other images while a run is going on cancels the run (\emph{The images changed: the
running analysis is cancelled.}), because its results would not belong to the new images.

\subsection{Cancel}

\ui{Cancel} is enabled while a run is going on. It stops the run at the next check point:
between two frames, between two tiles of a frame, or during a model download. The run does not
stop in the middle of a window, so it can take a moment.

Frames already computed are kept. The run returns them as its results, as a complete run would:
the \ui{LOG} reports \emph{stopped by the user after $k$ of $n$ pairs} and then \emph{Results
received: $k$ frame(s)}; \ui{Export Results} is enabled, the strain window opens, and the
tracking-quality check (Section~\ref{sec:tracking-quality}) runs on the computed frames. The
frames that were not computed have no result: the canvas shows their image without a field, the
strain window shows \emph{Frame $k$ was not computed.}, and exported data files hold NaN for them
(with a \file{Done} array that marks the computed frames, Section~\ref{sec:export}).

If you cancel before the first frame is finished, nothing replaces the previous results
(\emph{No frame was computed; the previous results are kept.}), and the strain window does not
open. Cancelling during the model
download stops the download; the \ui{LOG} says \emph{Stopped.}

\begin{warnbox}
A cancelled run that finished at least one frame replaces the previous results with its partial
ones. Save a project (\menu{File}{Save Project}, \kbd{Ctrl+S}) before you start a run you may
cancel if you still need the previous results.
\end{warnbox}

A run that stops on an error is different: it returns no frames at all. An \ui{Analysis failed}
dialog shows the cause, the \ui{LOG} holds the full traceback, the previous results, if any,
stay, and the strain window does not open.

\subsection{The PROGRESS area}

\begin{description}[leftmargin=1.2cm,style=nextline]
  \item[Progress bar]
    The fraction of the frame pairs done. It also moves inside a pair when the pair is cut into
    tiles.
  \item[Status line]
    \ui{Ready} when no run is going on. During a run it shows the percentage done and
    \ui{Frame $k$/$N$}: $k$ is the last frame finished and $N$ the number of images, frame~1
    included. Both count from 1, like the frame slider. Frame~1 is the reference, so the first
    frame to finish is frame~2 or later. Before the first frame finishes, the line shows the
    percentage alone.
  \item[Model download]
    While a run downloads its model (first use only), the status line shows \emph{Downloading
    model RAFTcorr-s4: $a$ of $b$ MB} instead of the frame, and the bar follows the download. A
    model that can be neither found nor downloaded stops the run and opens the Models window
    (Section~\ref{sec:installing-models}).
  \item[ELAPSED]
    The time since you clicked \ui{Run DIC Analysis}, in minutes and seconds, updated every
    second. It includes loading, and downloading, the model.
  \item[REMAINING]
    An estimate from the elapsed time and the fraction done; \ui{-{}-:-{}-} until 1\,\% of the
    run is done. Because the elapsed time includes loading the model, the first estimates are
    often too long; the estimate settles after a few frames.
\end{description}

\subsection{Large regions: tiles}
\label{sec:tiling}

The network sees a window of the reference image and the same window of the deformed image. A
point can only be matched if its displaced position lies inside that window, so the window is
the bounding box of the region of interest grown by a margin. By default the margin is the reach
of the model, 32\,$s$\,px (128\,px for RAFTcorr-s4); \ui{Margin} and \ui{Use the Model Reach}
under \ui{ADVANCED} change it (Section~\ref{sec:parameters}).

The memory one pass of the network needs grows with the number of pixels in the window, at a
rate that depends on the model: a larger $s$ needs less memory per pixel, and the stored
correlation volume that RAFTcorr-s8 uses by default grows faster than the window area. Before
each frame pair, RAFTcorr estimates the peak memory of one pass with a memory model of the
selected model and compares it with the share of the free GPU memory set by \ui{GPU Memory}
(default 0.85, under \ui{ADVANCED}), or with a fixed 6\,GB on the CPU.

\begin{itemize}
  \item If the whole window fits, the pair is solved in one window.
  \item If it does not, the region is cut into tiles. Each tile has a \emph{core}, the pixels it
        writes, and a \emph{halo} around it, the margin, which the network sees but whose result
        is not used. Neighbouring cores overlap by \ui{Tile Blending} (default 32\,px), and the
        overlap is blended with linear weights so that no seam shows.
\end{itemize}

Tiling is automatic and needs no setting. The \ui{LOG} says which case applies, at the first pair
and again whenever the plan changes (in incremental tracking the region of each reference frame
differs, and the free memory can change during a run):

\begin{itemize}
  \item \emph{tiling: one window of $W$ x $H$ px (margin $M$ px)};
  \item \emph{tiling: $n$ tiles, windows up to $W$ x $H$ px, margin $M$ px}.
\end{itemize}

\begin{warnbox}
When the memory only allows small windows, RAFTcorr reduces the halo so that each tile keeps a
useful core, and the tiling line ends with \emph{(margin reduced to fit memory)}. A point that
moves farther than the reduced margin can then leave the window of its tile and be matched
wrongly. Free GPU memory (close other GPU programs), raise \ui{GPU Memory}, choose a model with a
larger $s$, or reduce the region of interest.
\end{warnbox}

The memory model is an estimate. If the GPU still runs out of memory, RAFTcorr empties its memory
cache, makes the windows smaller (at most 55\,\% of the area of the largest window that failed) and
solves the pair again; the \ui{LOG} shows \emph{out of device memory; retrying with windows of at
most \ldots{} px}. After three such retries the run stops with an \ui{Analysis failed} dialog
that begins with \emph{The GPU ran out of memory. Close other GPU programs, choose a model with a
larger s, or lower GPU Memory in ADVANCED.} A budget so small that a window would fall below the
smallest size the model accepts also stops the run, with a message that gives both sizes and
begins with \emph{There is not enough free GPU memory for the smallest window this model accepts.
Close other GPU programs, or raise GPU Memory in ADVANCED if it is set low.}

\subsection{Tracking quality warning}
\label{sec:tracking-quality}

When a run ends, RAFTcorr checks the residual of every computed frame: the grey-level mismatch
that remains after the deformed image is mapped back onto frame~1 with the measured displacement
(Section~\ref{sec:residual}). It takes the median residual over the region of interest of each
frame. The residual grows slowly as the specimen deforms, but it jumps when tracking fails,
typically when the displacement from the reference frame exceeds the reach of the model.

A frame counts as a jump when its median residual exceeds 1.6 times that of the previous computed
frame and 2 times that of the first computed frame. From the first jump on, every frame whose
median residual stays above 1.6 times the level before the jump is flagged. The check needs at
least two computed frames after frame~1, and it runs after a cancelled run too, on the frames
computed.

When frames are flagged, a \ui{Tracking Quality Warning} dialog opens with the summary \emph{Part
of the sequence may not be tracked correctly.}, then the finding and an advice that depends on the
tracking mode:

\begin{itemize}
  \item \emph{The tracking quality drops from frame $k$ on: the intensity residual rises to $r$
        grey levels, against $r_0$ before.} Here $r_0$ is the highest median residual of the
        frames before frame~$k$.
  \item In accumulative tracking: \emph{The displacement from frame 1 may exceed the reach of the
        model. Try incremental tracking, or a model with a larger s.}
  \item In the other tracking modes: \emph{The step between reference frames may exceed the reach
        of the model. Use more reference frames, or a model with a larger s.}
\end{itemize}

\ui{Show Details\ldots} lists the flagged frames (the first ten, then how many more) and suggests
showing the \ui{Residual} field. Frame numbers count from 1, as on screen. The dialog does not
block the application: it sits in front of the strain window, you can keep working while it is
open, and the next run closes it. The \ui{LOG} keeps the same finding in one line
(\emph{tracking quality drops from frame $k$ on \ldots}), and the note is stored with the
provenance of the run in exported files.

The check never changes a result. When it fires:

\begin{enumerate}
  \item Select \ui{Residual} in \ui{FIELD} and step through the flagged frames.
  \item If the residual is high over most of the region, tracking was lost. In accumulative
        tracking, switch \ui{Tracking Mode} to \ui{Incremental}, or choose a model with a larger
        $s$ (\ui{Suggest a Model} measures the motion and picks one,
        Section~\ref{sec:workflow}). In incremental tracking, use more reference frames
        (\ui{Every Frame}, a smaller \ui{Interval}, or more \ui{Reference Frames}), or a model
        with a larger $s$.
  \item If the residual is high only where the images themselves change (a crack that opens,
        light reflected from the surface, speckle that flakes off), the displacement elsewhere
        can still be right. The median used by the check moves only when such an area is large;
        exclude it with a region of interest on the later frames (Section~\ref{sec:roi}).
\end{enumerate}

\subsection{When a run finishes}

At the end of a run that returns at least one computed frame:

\begin{itemize}
  \item the \ui{LOG} shows \emph{done: $n$ frames in \ldots{} s} and, in green, \emph{Results
        received: $n$ frame(s)}, where $n$ counts the computed frames, frame~1 excluded;
  \item \ui{Run DIC Analysis} is enabled again, \ui{Cancel} is disabled and \ui{Export Results}
        becomes available (also as \menu{File}{Export Results\ldots}, \kbd{Ctrl+E});
  \item the canvas moves to the last computed frame and draws the field selected in \ui{FIELD}
        (Section~\ref{sec:results});
  \item the strain window opens by itself, in front of the main window
        (Section~\ref{sec:strain}). It computes nothing until you click \ui{Compute Strain}. If
        you only need displacement, close it; \ui{Open Strain Window} brings it back;
  \item a strain result of the previous run is discarded, since it belongs to other
        displacements;
  \item the \ui{Tracking Quality Warning} dialog opens if frames were flagged.
\end{itemize}
